% Thesis Metadata
% Fill in your personal information and thesis details here

% Thesis type (uncomment one)
% \BSc  % For Bachelor of Science
 \MSc  % For Master of Science
% \PhD  % For Doctor of Philosophy

% Student Information
\author{Erlind}{Skura}  % {First Name}{Last Name}

% Thesis titles
\title{DEEP LEARNING FOR INSTANCE SEGMENTATION AND QUANTIFICATION OF BEAD DEFECTS IN SEM IMAGES OF ELECTROSPUN NANOFIBRE MATS}
\titleenglish{DEEP LEARNING FOR INSTANCE SEGMENTATION AND QUANTIFICATION OF BEAD DEFECTS IN SEM IMAGES OF ELECTROSPUN NANOFIBRE MATS}
\titlealbanian{MËSIMI I THELLË PËR SEGMENTIMIN E INSTANCAVE DHE KUANTIFIKIMIN E DEFEKTEVE TË TIPIT RRUAZË NË IMAZHE SEM TË MEMBRANAVE NANOFIBROZE TË ELEKTROTJERRURA}

% TODO(supervisor): title not yet approved.

% Dates
\date{September 2026}
\dateenglish{September 2026}

% Approval date (shown on the approval sheet next to "Date:") -- confirmed by Erlind.
\approvaldate{4 September 2026}

% Supervisor
\supervisor{Assoc. Prof. Dr. Arban Uka}
\supervisorenglish{Assoc. Prof. Dr. Arban Uka}

% Department and Faculty
\department{Computer Engineering}
\departmentenglish{Computer Engineering}
\departmentalbanian{Inxhinierisë Kompjuterike}
\faculty{Faculty of Architecture and Engineering}
\facultyenglish{Faculty of Architecture and Engineering}

% Head of Department (for approval page) -- confirmed by Erlind.
\headofdepartment{Dr. Florenc Skuka}

% Committee members (for approval page - Not required for BSc)
% TODO(supervisor): only the supervisor is known so far -- ask for the other two
% examiners once the defence is scheduled (question 20).
\committee
    {Assoc. Prof. Dr. Arban Uka}  % Chair/Supervisor
    {}
    {}
    {}
    {}

% Abstract (English)
\abstract{
Bead defects in electrospun nanofibre mats degrade porosity and mechanical
strength, and are quantified by hand, one ImageJ outline at a time.

This thesis automates that measurement. The dataset is 11 SEM micrographs of four
specimens at 500$\times$, 1000$\times$ and 3000$\times$, carrying 606 manually
delineated beads; since the specimens are the only independent units in it,
evaluation is leave-one-specimen-out. Mask R-CNN is compared under an identical
protocol against five alternatives: a detection-only ablation of itself, YOLOv8
and YOLOv5u, a U-Net semantic baseline and a classical Otsu-plus-watershed
pipeline.

Mask R-CNN leads the mask-predicting methods at
$\mathrm{AP}_{50} = 0.298 \pm 0.059$, a counting error of 41.1\% per micrograph
and a median bead diameter within 3.2\% of the manual value -- two thousandths
short of the 0.30 criterion, taken from a benchmark of matched object size. The
semantic baseline matches its mask accuracy but merges 48.7\% of beads and
under-counts by 44\%: what instance segmentation contributes is separation, not
boundary precision. Detection improves with magnification, 0.227 at $500\times$
against 0.396 at $3000\times$ although that group holds only 52 beads: the
binding constraint is pixels per object, not label scarcity, and magnifying the
dataset threefold raises it to 0.373 without adding an observation. YOLOv5u, the
best detector tested, counts three specimens to within 18.3\% and the fourth to
212.5\%.
}

% Keywords (English)
\keywordsenglish{Instance Segmentation, Electrospun Nanofibres, Bead Defects, Scanning Electron Microscopy, Mask R-CNN}

% Abstrakt (Albanian - optional, Not required for BSc)
\abstrakt{
Defektet e tipit rruazë në membranat nanofibroze ulin porozitetin dhe
qëndrueshmërinë mekanike dhe maten manualisht, duke vizatuar secilën në ImageJ.

Kjo tezë e automatizon atë matje. Grupi ka 11 imazhe SEM të katër mostrave në
500$\times$, 1000$\times$ dhe 3000$\times$, me 606 rruaza të anotuara
manualisht; meqë mostrat janë njësitë e vetme të pavarura, vlerësimi lë jashtë
një mostër të plotë. Mask R-CNN krahasohet nën të njëjtin protokoll me pesë
alternativa: vetë ai pa degën e maskave, YOLOv8, YOLOv5u, një model bazë U-Net
dhe një metodë klasike Otsu-\emph{watershed}.

Mask R-CNN kryeson mes metodave me maska, me
$\mathrm{AP}_{50} = 0{,}298 \pm 0{,}059$, gabim numërimi 41,1\% për imazh dhe
diametër mesatar brenda 3,2\% të matjes manuale -- dy të mijta nën kriterin 0,30,
marrë nga një bazë të dhënash me objekte po aq të vogla. Modeli bazë semantik
arrin të njëjtën saktësi maskash, por bashkon 48,7\% të rruazave: kontributi i
segmentimit të instancave është ndarja e objekteve, jo konturi. Detektimi
përmirësohet me zmadhimin, 0,227 në $500\times$ kundrejt 0,396 në $3000\times$
edhe pse ai grup ka vetëm 52 rruaza: kufizimi është numri i pikselëve për objekt,
jo etiketat e pakta; zmadhimi i grupit me faktorin tre e ngre në 0,373. YOLOv5u,
detektori më i mirë i testuar, i numëron tri mostra brenda 18,3\% dhe të katërtën
me 212,5\%.
}

% Keywords (Albanian)
\keywords{Segmentim Instancash, Nanofibra të Elektrotjerrura, Defekte Rruazë, Mikroskopi Elektronike me Skanim, Mask R-CNN}

% Acknowledgments
\acknowledgments{
I would like to thank my supervisor, Assoc. Prof. Dr. Arban Uka, for proposing this
problem, for providing the micrographs and the annotations on which the whole study
rests, and for his guidance throughout the work.

I am also grateful to my family for their support during my studies, and to Epoka
University for the resources that made this research possible.
}

% Abbreviations
\abbreviations{
\begin{longtable}{>{\bfseries}ll}
AJI & Aggregated Jaccard Index \\
AP & Average Precision \\
CNN & Convolutional Neural Network \\
EHT & Extra High Tension (accelerating voltage) \\
FPN & Feature Pyramid Network \\
GPU & Graphics Processing Unit \\
IoU & Intersection over Union \\
LOSO & Leave-One-Specimen-Out \\
mAP & Mean Average Precision \\
NMS & Non-Maximum Suppression \\
PQ & Panoptic Quality \\
R-CNN & Region-based Convolutional Neural Network \\
RoI & Region of Interest \\
RPN & Region Proposal Network \\
SAM & Segment Anything Model \\
SEM & Scanning Electron Microscopy \\
WD & Working Distance \\
\end{longtable}
}

% Optional: Disable table or figure lists if not needed
% \NoTableList
% \NoFigureList
